\begin{block}{Morphology-Based Kugire Annotation}
Searches the morphological data for positions that may form a kugire. Conclusive
forms, imperative forms, realis forms, sentence-final particles, etc.\ serve as
strong candidates; binding particles, kakari-musubi, taigen-dome, etc.\ serve as
moderate-to-weak candidates. This method emphasizes ``grammatically observable
break points.''

\vspace{1.0ex}

\begin{figure}[H]\centering
\includesvg[inkscapelatex=false, width=0.9\textwidth]{../images/example}
\setcounter{figure}{6}
\caption{Three classes of kugire candidates derived from morphological information:
strong, moderate, and weak candidates. This figure illustrates that the rule-based
method does not give a single answer, but candidate interpretations with certainty
levels.}
\end{figure}
\end{block}
